%!TEX program = lualatex
% Rascunho: preencher as duas massas dos filtros antes de concluir os resultados.
\documentclass[a4paper,12pt]{article}
\usepackage[brazilian]{babel}
\usepackage{fontspec}
\setmainfont{Latin Modern Roman}
\usepackage{microtype}
\usepackage[top=3cm,bottom=2.5cm,left=3cm,right=2.5cm]{geometry}
\usepackage{setspace,indentfirst}
\onehalfspacing
\usepackage{amsmath,amssymb}
\usepackage[version=4]{mhchem}
\usepackage{siunitx}
\sisetup{locale=DE,output-decimal-marker={,},per-mode=symbol}
\usepackage{booktabs,array}
\usepackage{graphicx}
\usepackage{hyperref}
\hypersetup{hidelinks,pdftitle={Prática 9 -- Determinação gravimétrica de sulfato em água do mar}}

\begin{document}
\begin{titlepage}
\centering
{\large Universidade do Estado do Rio de Janeiro -- UERJ\par}
\vspace{0.3cm}{\large Química VIII / Oceanografia\par}
\vspace{1.0cm}\includegraphics[width=0.28\linewidth]{uerj.jpg}\par
\vspace{0.8cm}{\LARGE\bfseries Prática 9\par}
\vspace{0.5cm}{\Large\bfseries DETERMINAÇÃO GRAVIMÉTRICA DE SULFATO\\ EM ÁGUA DO MAR\par}
\vfill
\begin{flushright}
\textbf{Alunos:}\\ Israel Carneiro\\ José Mauro Xavier Elsas\\ Miguel Gayani\\[0.5cm]
\textbf{Professora:} Priscilla Magalhães de Souza
\end{flushright}
\vspace{1cm}Rio de Janeiro -- 2026
\end{titlepage}

\section{Introdução}
A água do mar contém principalmente água e sais dissolvidos; entre os íons
mais abundantes estão \ce{Cl-}, \ce{Na+}, \ce{SO4^{2-}}, \ce{Mg^{2+}},
\ce{Ca^{2+}} e \ce{K+}. A proporção desses constituintes pode variar com
mistura, evaporação e processos biogeoquímicos. O sulfato é um dos
constituintes dissolvidos da água do mar. Sua determinação
gravimétrica baseia-se na formação de sulfato de bário, pouco solúvel, a partir
da adição de íons bário à amostra:
\begin{equation}
  \ce{Ba^{2+}(aq) + SO4^{2-}(aq) -> BaSO4(s)}.
\end{equation}
A relação em quantidade de matéria entre o sulfato presente na alíquota e o
\ce{BaSO4} formado é de 1:1. A precipitação a quente e a digestão favorecem
a obtenção de partículas filtráveis. O meio ácido ajuda a reduzir a
interferência de carbonatos, conforme discutido nas anotações da aula:
\ce{CO3^{2-} + H+ -> HCO3-} e
\ce{HCO3- + H+ -> H2CO3 -> CO2 + H2O}.

Uma vantagem da gravimetria é medir diretamente a massa de um composto de
composição conhecida, relacionando-a à massa do analito por estequiometria,
sem construir uma curva de calibração instrumental. Para que essa vantagem
se concretize, a precipitação deve ser completa e o sólido deve estar livre
de contaminantes e com massa estabilizada após secagem.

Também se pode determinar sulfato por cromatografia de íons (separação e
quantificação do ânion), por turbidimetria (medida da turbidez do
\ce{BaSO4} disperso) e por métodos complexométricos indiretos. A apostila
apresenta uma prática específica de dosagem complexométrica de sulfato;
a EPA descreve métodos de cromatografia de íons e turbidimetria nas
referências finais.

\section{Objetivos}
\begin{itemize}
  \item precipitar o sulfato presente em uma amostra de água do mar na forma de
        \ce{BaSO4};
  \item calcular a massa de precipitado após descontar a massa de cada filtro;
  \item determinar a concentração de \ce{SO4^{2-}} em
        \si{\milli\gram\per\litre} e em porcentagem massa/volume.
\end{itemize}

\section{Material utilizado}
\begin{itemize}
  \item béquer de \SI{400}{\milli\litre} e béquer de
        \SI{50}{\milli\litre}, indicados no quadro;
  \item material para medir e transferir os volumes da amostra e das soluções;
  \item utensílio para agitação durante a adição do precipitante;
  \item equipamento de aquecimento;
  \item dispositivo de filtração e dois filtros de massas previamente medidas;
  \item equipamento de secagem e balança para a pesagem dos filtros e
        precipitados.
\end{itemize}
O quadro não especifica os modelos da vidraria volumétrica, do dispositivo
de filtração e do equipamento de secagem. A apostila apresenta uma relação
mais detalhada de materiais, mas descreve uma versão diferente do
procedimento; por isso, esses detalhes não foram tomados como confirmação
do que foi utilizado em sala.

\section{Reagentes e soluções}
\begin{itemize}
  \item amostra de água do mar;
  \item água para diluição da amostra e lavagem do precipitado;
  \item solução de \ce{HCl} \SI{3}{\mole\per\litre};
  \item solução de \ce{BaCl2} 4\% m/v.
\end{itemize}
O registro da aula não distingue a forma anidra da hidratada do cloreto de
bário empregado no preparo da solução.

\section{Procedimento experimental}
O presente relatório corresponde ao grupo da capela 2. O quadro e os dados
compartilhados registram pesagens das capelas 1 e 2; ambas foram mantidas
para o cálculo da média solicitado, sujeita à confirmação de que as duas
capelas analisaram alíquotas da mesma amostra.

\subsection{Precipitação do sulfato}
O roteiro do quadro da aula, datado de 21/09/2026, indica adicionar a um
béquer de \SI{400}{\milli\litre} uma alíquota de \SI{25}{\milli\litre} de
água do mar, \SI{100}{\milli\litre} de água e \SI{2}{\milli\litre} de
\ce{HCl} \SI{3}{\mole\per\litre}. Aquecer sem deixar entrar em ebulição.
Adicionar aproximadamente \SI{25}{\milli\litre} de solução de \ce{BaCl2}
4\% m/v, gota a gota sob agitação. Verificar a precipitação completa
pela adição de mais precipitante ao sobrenadante, observando a eventual
formação de sólido.

\subsection{Digestão, filtração e pesagem}
Deixar o precipitado em digestão por cerca de 30 minutos, filtrar e secar.
Pesar os filtros vazios e, após a separação e a secagem, pesar cada conjunto
filtro + precipitado.
O registro fornecido informa duas pesagens finais de filtro com precipitado
e as massas dos filtros vazios. A condição exata de secagem não foi informada.

\noindent\textit{Critério de registro:} a apostila identifica o experimento
como Prática 11 e descreve outras quantidades (por exemplo, alíquota de
\SI{50.00}{\milli\litre}). Este relatório o denomina Prática 9 e usa as
quantidades do quadro da aula. Detalhes de operação presentes apenas na
apostila não foram atribuídos à execução da turma.

\section{Resultados e cálculos}
As duas massas iniciais correspondem ao conjunto filtro + precipitado.
As massas dos filtros vazios foram informadas posteriormente, em gramas.
\begin{table}[h]
\centering
\caption{Pesagens e correção pela massa do filtro.}
\begin{tabular}{c c c c}
\toprule
Capela & Filtro + \ce{BaSO4} (g) & Filtro vazio (g) & \ce{BaSO4} (g)\\
\midrule
1 & 0,7337 & 0,5990 & 0,1347\\
2 & 0,7515 & 0,6150 & 0,1365\\
\bottomrule
\end{tabular}
\end{table}

As massas dos filtros foram descontadas \emph{antes} de calcular a média:
\begin{equation}
  \overline m_{\ce{BaSO4}} =
  \frac{(0,7337-0,5990)+(0,7515-0,6150)}{2}
  =\frac{0,1347+0,1365}{2}=\SI{0.1356}{\gram}.
\end{equation}
As massas individuais de precipitado diferem em \SI{0.0018}{\gram},
aproximadamente 1,33\% da média.

Empregando as massas molares registradas no quadro,
$M(\ce{SO4^{2-}})=\SI{96.06}{\gram\per\mole}$ e
$M(\ce{BaSO4})=\SI{233.39}{\gram\per\mole}$,
\begin{equation}
 m_{\ce{SO4^{2-}}}
 =\SI{0.1356}{\gram}\,
 \frac{\SI{96.06}{\gram\per\mole}}{\SI{233.39}{\gram\per\mole}}
 =\SI{0.05581}{\gram}.
\end{equation}
Para uma alíquota de \SI{25}{\milli\litre}, a concentração na amostra
original, expressa em mg/L, é
\begin{equation}
 C_{\ce{SO4^{2-}}}=\frac{\SI{0.1356}{\gram}}{\SI{0.025}{\litre}}
 \frac{96,06}{233,39}\times 1000\ \mathrm{mg\,g^{-1}}
 \simeq\SI{2232}{\milli\gram\per\litre}.
 \label{eq:conc}
\end{equation}
Em porcentagem massa/volume (gramas de sulfato por 100 mL da amostra):
\begin{equation}
 \%\,(m/v)=0,1356\frac{96,06}{233,39}\frac{100}{25}
 \simeq 0,2232\%.
\end{equation}
Individualmente, as capelas 1 e 2 correspondem, respectivamente, a
\SI{2218}{\milli\gram\per\litre} e
\SI{2247}{\milli\gram\per\litre}, considerando alíquotas de
\SI{25}{\milli\litre} da mesma amostra.
O volume de água adicionado serve à precipitação e não substitui os
\SI{25}{\milli\litre} da amostra no denominador.

\subsection*{Preparo das soluções: relações solicitadas no quadro}
Uma solução de \ce{BaCl2} 4\% m/v contém \SI{4}{\gram} de soluto em
\SI{100}{\milli\litre} de solução, ou \SI{40}{\gram\per\litre}.
Se o reagente for \ce{BaCl2} anidro ($M=\SI{208.23}{\gram\per\mole}$),
$C=40/208,23\simeq\SI{0.1921}{\mole\per\litre}$.
Se for \ce{BaCl2.2H2O} ($M\simeq\SI{244.26}{\gram\per\mole}$),
$C\simeq\SI{0.1638}{\mole\per\litre}$. Conferir o rótulo antes de
escolher uma das duas conversões.
Para preparar \ce{HCl} \SI{3}{\mole\per\litre} por diluição,
$C_1V_1=C_2V_2$ e $V_1=(3,00\,V_2)/C_1$; a concentração de partida e o
volume final preparado ainda não foram informados.

\section{Discussão}
O método vincula diretamente a massa corrigida de \ce{BaSO4} à quantidade de
sulfato da alíquota. As duas massas corrigidas, \SI{0.1347}{\gram} e
\SI{0.1365}{\gram}, produziram uma média de \SI{0.1356}{\gram} e uma
concentração média estimada de \SI{2232}{\milli\gram\per\litre} de
\ce{SO4^{2-}} na amostra original. O quadro pede conferência da
precipitação completa com excesso de \ce{BaCl2}; perdas na transferência
e retenção de água ou impurezas no filtro também alterariam a massa final.
O uso da média pressupõe que as duas capelas analisaram alíquotas da mesma
amostra sob as condições registradas no quadro. Essa informação deve ser
confirmada com o grupo antes da entrega definitiva.

\section{Conclusão}
Após a correção pelas massas dos filtros, a massa média de \ce{BaSO4} foi
\SI{0.1356}{\gram}. Considerando alíquotas de \SI{25}{\milli\litre}
da mesma amostra, a concentração calculada de sulfato na água do mar foi
aproximadamente \SI{2232}{\milli\gram\per\litre}, equivalente a
0,2232\% m/v. A identificação comum da amostra nas duas capelas ainda
precisa ser confirmada.

\section*{Referência}
LAGO, Dalva Cristina Baptista do; GOIS, Jefferson Santos de; LANGONE,
Marta Antunes Pereira; SENNA, Lilian Ferreira de.
\textit{Química VIII -- Oceanografia: aulas práticas}. Rio de Janeiro:
UERJ, Instituto de Química, Departamento de Química Analítica, 2017,
Prática 11, pp.~47--49.

\noindent NATIONAL OCEANIC AND ATMOSPHERIC ADMINISTRATION (NOAA).
\textit{Concentration of the major components of seawater}, tabela 1.1.
Disponível em: \url{https://repository.library.noaa.gov/view/noaa/12066/noaa_12066_DS1.pdf}.

\noindent UNITED STATES ENVIRONMENTAL PROTECTION AGENCY (EPA).
\textit{Method 300.0: Determination of inorganic anions by ion chromatography},
rev. 2.1, 1993. Disponível em:
\url{https://www.epa.gov/sites/default/files/2015-08/documents/method_300-0_rev_2-1_1993.pdf}.

\noindent UNITED STATES ENVIRONMENTAL PROTECTION AGENCY (EPA).
\textit{SW-846 Test Method 9038: Sulfate (Turbidimetric)}.
Disponível em:
\url{https://www.epa.gov/hw-sw846/sw-846-test-method-9038-sulfate-turbidimetric}.
\end{document}
